\Def Поле $K$ называется алгебраически замкнутым, если выполнено одно из следующих эквивалентных условий:

\begin{enumerate}
\item Любой многочлен $f(x) \in K[x]$ с $deg f (x) \geqslant 1$ раскладывается на линейные множители в $K$.
\item  Любой многочлен $f(x) \in K[x]$ с $deg f (x) \geqslant 1$ имеет корень в $K$.
\item Все неприводимые над $K$ многочлены имеют степень 1.
\item Любое алгебраическое расширение над $K$ тривиально (то есть совпадает с $K$).
\item Любое конечное расширение над $K$ тривиально.
\item Для любого многочлена $f(x) \in K[x]$ с $deg f (x) \geqslant 1$ его поле разложения совпадает с $K$.
\end{enumerate}

\Def Алгебраическим замыканием поля $F$ называется алгебраическое расширение $K = \overline{F}$ поля $F$, которое является алгебраически замкнутым полем.

\Proof 
$(1) \Rightarrow (2)$: Тривиально, если многочлен $f: deg f(x) \geqslant 1$ раскладывается на линейные сомножители, то он имеет корень.

$(2) \Rightarrow (3)$: Если многочлен $f$ имеет корень $\alpha \in K$, то $x - \alpha | f$. Тогда, в силу неприводимости, $f \simeq x - \alpha$.

$(3) \Rightarrow (4)$: Пусть $L \supset K$ — алгебраическое расширение, $\alpha \in L$. Рассмотрим минимальный многочлен $m_{\alpha} \in K[x]$ элемента $\alpha$ над $K$. Поскольку $m_{\alpha}$ неприводим, то $deg m_{\alpha} = 1$, откуда $\alpha \in K$.

$(4) \Rightarrow (5)$: Тривиально, поскольку любое конечное расширение является алгебраическим.

$(5) \Rightarrow (6)$: Поле разложения многочлена $f$ является конечным расширением поля $K$, а значит, совпадает с $K$.

$(6) \Rightarrow (1)$: Поле разложения многочлена $f$ совпадает с $K$, значит, $f$ раскладывается на линейные сомножители над $K[x]$.
\EndProof